
Comme les deux membres de la relation à démontrer sont additifs, on peut supposer que @@M0@@ est la classe d'un @@M1@@-Module localement libre de type fini @@M2@@. L'élément @@M3@@ est alors la classe du complexe @@M4@@, obtenu en prenant une résolution projective du complexe @@M5@@, et en la tensorisant par @@M6@@. Les faisceaux d'homologie de ce complexe sont donc les @@M7@@. D'après 2.4, on a l'isomorphisme:
@@M8@@
D'après 2.5, on a l'isomorphisme:
@@M9@@

Les @@M10@@ sont donc des @@M11@@-Modules localement libres de type fini, et la classe dans @@M12@@ du complexe @@M13@@ est la somme alternée des classes de ses faisceaux d'homologie, soit:
@@M14@@

\textbf{Corollaire 2.8.} \uline{Sous les hypothèses de} 2.7, \uline{soient} @@M15@@; \uline{on a la relation}:
@@M16@@

D'après 2.7, @@M17@@.

Par application de la formule de projection, on trouve donc:
@@M18@@

\section*{3. Calcul du \texorpdfstring{@@M19@@}{K} d'un schéma éclaté}

\textbf{3.1.} Dans ce paragraphe, on désigne par @@M20@@ une immersion fermée régulière. On note @@M21@@ le schéma obtenu en faisant éclater @@M22@@, et @@M23@@.
